\documentclass[12pt]{article}
\usepackage[T1]{fontenc}
\usepackage[utf8]{inputenc}
\usepackage[brazil]{babel}
\usepackage[margin=1in]{geometry}
\usepackage{newtxtext,newtxmath}
\usepackage{ragged2e}
\usepackage[normalem]{ulem}
\setlength{\parindent}{0pt}
\linespread{1.15}
\setlength{\parskip}{0.7\baselineskip}
\emergencystretch=3em
\sloppy
\begin{document}
\justifying
\noindent \textbf{Processo n.: }\uline{5003714-34.2025.4.02.5101}

\noindent Verifica-se que a questão objeto da controvérsia jurídica suscitada no recurso manejado pela parte recorrente esteve afetada no representativo da controvérsia - \textbf{Tema n. }\textbf{309} da TNU - (PU no processo n. PEDILEF 5001816-07.2020.4.04.7008/PR),  afetado em 18/08/2022 foi julgado (trânsito em julgado em 14/02/2024), por meio do qual foi elucidada a seguinte questão:

\noindent \textit{"O auxílio-alimentação integra a base de cálculo da licença-prêmio não usufruída e convertida em pecúnia?"}

\noindent Restou firmada a seguinte tese:

\noindent \textit{"O auxílio-alimentação pago aos servidores públicos federais (Lei n. 8.460/92) integra a base de cálculo da licença-prêmio não usufruída e convertida em pecúnia." (Tema 309 TNU).}

\noindent Com efeito, depreende-se da leitura do voto condutor do acórdão que o mesmo se encontra em conformidade com o decidido no  tema pela TNU.

\noindent A proposito, consta do acordao hostilizado: (...) Auxílio-alimentação, Sistema Remuneratório e Benefícios, Servidor Público Civil, DIREITO ADMINISTRATIVO E OUTRAS MATÉRIAS DE DIREITO PÚBLICO

\noindent Nessas condições, o conteúdo do Acórdão recorrido é consentâneo com o entendimento da TNU, firmado em julgamento de recurso representativo de controvérsia, o que atrai a incidência da regra prevista pelo art. 14, III, \textit{b,} da Resolução CJF n. 586/2019:

\noindent \textit{\textbf{Art. 14.}}\textit{ Decorrido o prazo para contrarrazões, os autos serão conclusos ao magistrado responsável pelo exame preliminar de admissibilidade, que deverá, de forma sucessiva: (...) }\textit{\textbf{III }}\textit{- negar seguimento a pedido de uniformização de interpretação de lei federal interposto contra acórdão que esteja em conformidade com entendimento consolidado: }\textit{\textbf{b)}}\textit{ em }\textit{\textbf{\uline{recurso representativo de controvérsia pela Turma Nacional de Uniformização}}}\textit{ ou em pedido de uniformização de interpretação de lei dirigido ao Superior Tribunal de Justiça;(...).}

\noindent Posto isso, com arrimo no \textbf{\uline{art. 14, III, }}\textit{\textbf{\uline{b}}}\textit{,} da Resolução CJF n. 586/2019 c/c art. 1.030, I, do CPC, \textbf{NEGO SEGUIMENTO }\textbf{a}o\textbf{ PU}\textbf{.}

\noindent Intimem-se. Aguarde-se o decurso do prazo recursal. Após, certifique-se o trânsito em julgado e devolva-se ao juízo de origem.
\end{document}
